\documentclass[10pt,a4paper]{article}
\usepackage[margin=18mm,headheight=14pt]{geometry}
\usepackage{fontspec}
\setmainfont{lmroman10-regular.otf}[BoldFont=lmroman10-bold.otf,ItalicFont=lmroman10-italic.otf,BoldItalicFont=lmroman10-bolditalic.otf]
\setsansfont{lmsans10-regular.otf}[BoldFont=lmsans10-bold.otf]
\setmonofont{lmmonolt10-regular.otf}[BoldFont=lmmonolt10-bold.otf,ItalicFont=lmmonolt10-oblique.otf]
\usepackage[ngerman]{babel}
\usepackage{amsmath,amssymb,booktabs,tabularx,array,fancyhdr,listings,xcolor}
\usepackage[hidelinks]{hyperref}
\pagestyle{fancy}\fancyhf{}
\fancyhead[L]{\sffamily AI-Hardware · Tag 4}
\fancyhead[R]{\sffamily Zahlensysteme und Binärarithmetik}
\fancyfoot[L]{\small Nachschlagen · Verstehen · In Python prüfen}
\fancyfoot[R]{\thepage\ / 6}
\setlength{\parindent}{0pt}\setlength{\parskip}{5pt}
\setlength{\tabcolsep}{7pt}\renewcommand{\arraystretch}{1.18}
\lstset{language=Python,basicstyle=\ttfamily\small,
 keywordstyle=\bfseries,commentstyle=\itshape\color{black},
 stringstyle=\ttfamily\color{black},showstringspaces=false,
 frame=single,rulecolor=\color{black},framesep=5pt,
 columns=fullflexible,keepspaces=true,breaklines=true,tabsize=4,
 aboveskip=8pt,belowskip=8pt}
\newcommand{\titlepagepart}[2]{\vspace*{-4mm}{\sffamily\small #1}\par{\sffamily\LARGE\bfseries #2}\par\vspace{2mm}}
\newcommand{\block}[1]{\par\vspace{2mm}{\sffamily\large\bfseries #1}\par}
\newcommand{\code}[1]{\texttt{#1}}
\begin{document}
\titlepagepart{01 / GRUNDLAGEN}{Zahlen schreiben und Bits lesen}
\textbf{Wozu?} Ein Register speichert Bits. Du brauchst eine Bitbreite und eine
Interpretation, um daraus eine Zahl zu lesen. Binär und Hex sind Schreibweisen;
unsigned und Zweierkomplement bestimmen die Bedeutung des Bitmusters.

\block{Stellenwerte: Nur die Stellen mit einer 1 mitzählen}
\begin{center}
\begin{tabular}{lrrrrrrrr}\toprule
Position &7&6&5&4&3&2&1&0\\
Unsigned &128&64&32&16&8&4&2&1\\
Beispiel &0&1&0&1&1&0&1&1\\\bottomrule
\end{tabular}
\end{center}
\[0101\,1011_2=64+16+8+2+1=91_{10}\]
Die rechte Stelle zählt $2^0=1$. Jede Stelle nach links zählt doppelt so viel.
Dezimal funktioniert genauso, nur mit $1,10,100,\ldots$ statt $1,2,4,\ldots$.

\block{Dezimal nach Binär}
Zerlege die Zahl in Zweierpotenzen: $45=32+8+4+1$.
Mit den Stellen $32,16,8,4,2,1$ ergibt das $101101_2$; auf acht Bits aufgefüllt:
\code{0010 1101}. Alternativ: wiederholt durch 2 teilen, Reste von unten nach oben lesen.

\block{Hex: Eine Ziffer entspricht vier Bits}
\begin{center}
\begin{tabular}{ccc|ccc|ccc|ccc}\toprule
Hex&Dez.&Bits&Hex&Dez.&Bits&Hex&Dez.&Bits&Hex&Dez.&Bits\\\midrule
0&0&0000&4&4&0100&8&8&1000&C&12&1100\\
1&1&0001&5&5&0101&9&9&1001&D&13&1101\\
2&2&0010&6&6&0110&A&10&1010&E&14&1110\\
3&3&0011&7&7&0111&B&11&1011&F&15&1111\\\bottomrule
\end{tabular}
\end{center}
Von rechts in Vierergruppen teilen; eine unvollständige Gruppe links mit Nullen ergänzen.
\[\texttt{0xB7}=1011\,0111_2=11\cdot16+7=183_{10}\]
\textbf{Achtung:} \code{0x} ist das Hexpräfix, \code{0b} das Binärpräfix.
\code{0x23} ist $2\cdot16+3=35$, nicht 23 dezimal. Hex legt kein Vorzeichen fest.

\block{Breite und Einheiten}
$N$ Bits haben $2^N$ Muster. Unsigned reichen sie von $0$ bis $2^N-1$;
die Null zählt als eine Möglichkeit mit.
\begin{center}
\begin{tabular}{ll|ll}\toprule
1 Nibble &4 Bits = 1 Hexziffer&1 Byte&8 Bits = 2 Hexziffern\\
LSB&niederwertigstes Bit rechts&MSB&höchstwertigstes Bit links\\
1 KiB&$1024=2^{10}$ Byte&1 kB&1000 Byte\\\bottomrule
\end{tabular}
\end{center}
Ein Word ist eine architekturabhängige Dateneinheit. $1\,\mathrm{MiB}=2^{20}$ Byte,
$1\,\mathrm{GiB}=2^{30}$ Byte. Kleinbuchstabe b steht häufig für Bit, B für Byte.

\newpage
\titlepagepart{02 / INTERPRETATION}{Negative Zahlen und das Warum}
\block{Warum Zweierkomplement?}
Wir wollen dieselbe Binäraddition auch für negative Zahlen verwenden.
Mit vier Bits wirkt das Muster \code{1111} beim Addieren wie „eins abziehen“:
\[
\underbrace{0101_2}_{5}+\underbrace{1111_2}_{\text{Muster für }-1}
=1\,0100_2\quad\longrightarrow\quad0100_2=4.
\]
Nur die rechten vier Bits werden gespeichert. Unsigned wäre dies $5+15=20$;
nach Wegfall der zusätzlichen 16 bleibt 4. Deshalb ist 15 als Bitmuster
praktisch für $-1$. Allgemein repräsentiert $2^N-k$ die Zahl $-k$, wenn sie
im signed Bereich liegt.

\block{Die linke Stelle zählt negativ}
\begin{center}
\begin{tabular}{lrrrrl}\toprule
Interpretation&\multicolumn{4}{c}{Stellenwerte bei vier Bits}&Muster \code{1101}\\\midrule
Unsigned&$+8$&4&2&1&$8+4+1=13$\\
Zweierkomplement&$-8$&4&2&1&$-8+4+1=-3$\\\bottomrule
\end{tabular}
\end{center}
\textbf{Kein bloßes Minuszeichen:} Im Zweierkomplement sind die restlichen Bits
nicht der Betrag. \code{1101} bedeutet $-3$, nicht $-5$.
Das System „Vorzeichen und Betrag“ würde $-5$ ergeben und hat zwei Nullen;
es ist nicht das System unserer Funktionen.

\block{Die Bitbreite gehört zur Bedeutung}
\begin{center}
\begin{tabular}{lrrr}\toprule
Muster&Bits&Unsigned&Zweierkomplement\\\midrule
\code{1101}&4&13&$-3$\\
\code{0000 1101}&8&13&$+13$\\
\code{1111 1101}&8&253&$-3$\\\bottomrule
\end{tabular}
\end{center}
Mit $N$ Bits ist der höchste Stellenwert $-2^{N-1}$. Daraus folgen:
\[
\boxed{\text{unsigned: }0\ldots 2^N-1}\qquad
\boxed{\text{signed: }-2^{N-1}\ldots 2^{N-1}-1}
\]
\begin{center}
\begin{tabular}{rll}\toprule
Bits&Unsigned&Zweierkomplement\\\midrule
1&0 bis 1&$-1$ bis 0\\
4&0 bis 15&$-8$ bis 7\\
5&0 bis 31&$-16$ bis 15\\
8&0 bis 255&$-128$ bis 127\\\bottomrule
\end{tabular}
\end{center}
\block{Negieren und Verbreitern}
\textbf{Vorzeichen umkehren:} Bits invertieren, 1 addieren, $N$ Bits behalten.
\[+3:\;0011\ \longrightarrow\ 1100\ \longrightarrow\ 1101:\;-3\]
Ausnahme: Die kleinste negative Zahl hat bei gleicher Breite kein positives
Gegenstück. Vier-Bit-$-8$ lässt sich nicht als Vier-Bit-$+8$ negieren.

\textbf{Vorzeichenerweiterung:} Links das bisherige Vorzeichenbit wiederholen.
\code{1101} wird \code{1111 1101} ($-3$ bleibt $-3$).
Positive Zahlen bekommen links Nullen. Unsigned Zahlen werden mit Nullen erweitert.

\newpage
\titlepagepart{03 / RECHNEN}{Addition, Carry und signed Overflow}
\block{Binäre Addition und Subtraktion}
$0+0=0$, $0+1=1$, $1+1=10_2$.
Bei $1+1$ schreiben wir 0 und tragen 1 nach links. Mit eingehendem Übertrag:
$1+1+1=11_2$; 1 schreiben und 1 weitertragen.

Subtraktion wird Addition der negierten zweiten Zahl: $a-b=a+(-b)$.
\[5-3:\quad0101+1101=1\,0010\quad\longrightarrow\quad0010=2.\]
\block{Zwei unterschiedliche Fragen}
\textbf{Carry-out:} Wird bei Addition der $N$-Bit-Muster eine 1 aus der höchsten
Stelle herausgetragen? Bei unsigned Addition zeigt dies einen Überlauf an.
Ein Übertrag \emph{innerhalb} der Rechnung ist noch kein Carry-out.

\textbf{Signed Overflow:} Liegt die mathematisch richtige Summe außerhalb des
signed Zahlenbereichs? Zwei positive Operanden können dann ein negatives Ergebnis
liefern; zwei negative ein nichtnegatives. Unterschiedliche Vorzeichen führen
bei gültigen Operanden nicht zu signed Overflow.
\begin{center}
\begin{tabular}{lclcc}\toprule
Signed Rechnung&Gespeichert&Wert&Carry-out&Signed Overflow\\\midrule
$-1+1$&\code{0000}&0&Ja&Nein\\
$7+1$&\code{1000}&$-8$&Nein&Ja\\
$-8+(-1)$&\code{0111}&7&Ja&Ja\\
$2+(-5)$&\code{1101}&$-3$&Nein&Nein\\\bottomrule
\end{tabular}
\end{center}
Alle Beispiele verwenden vier Bits. Die \textbf{richtige Summe} und der
\textbf{gespeicherte signed Wert} sind bei Overflow verschieden.

\block{Was die Python-Funktion nachbildet}
Python-Integer sind nicht von selbst auf unsere Registerbreite begrenzt.
Für die mathematische Summe $s=a+b$ gilt:
\[
\begin{aligned}
\mathrm{overflow}&=(s<-2^{N-1})\ \lor\ (s>2^{N-1}-1),\\
\mathrm{raw}&=s\bmod 2^N,\\
\mathrm{ergebnis}&=\operatorname{from\_twos\_complement}(\mathrm{raw},N).
\end{aligned}
\]
\code{raw} beschreibt das gespeicherte Muster als unsigned Zahl. Erst die letzte
Zeile interpretiert es als signed Zahl.

\block{Warum Modulo (\texttt{\%})?}
Bei vier Bits entspricht „die vier rechten Bits behalten“ dem Rest modulo 16:
\begin{center}
\begin{tabular}{rlrl}\toprule
Summe&Rechnung&\code{raw}&Signed Interpretation\\\midrule
18&$18=1\cdot16+2$&2&\code{0010} $\rightarrow 2$\\
8&$8=0\cdot16+8$&8&\code{1000} $\rightarrow -8$\\
$-3$&$-3=(-1)\cdot16+13$&13&\code{1101} $\rightarrow -3$\\
$-9$&$-9=(-1)\cdot16+7$&7&\code{0111} $\rightarrow 7$\\\bottomrule
\end{tabular}
\end{center}
Bei positivem Divisor 16 liefert Python mit \code{\% 16} immer einen Rest von
0 bis 15. Auch bei negativen Summen entsteht so das richtige gespeicherte Muster.

\block{Ein vollständiger Ablauf}
\code{add\_with\_overflow\_flag(7, 1, 4)}:
Summe $8$; Bereich $-8\ldots7$; Overflow wahr; $8\bmod16=8$;
Muster \code{1000}; signed Wert $-8$; Rückgabe \code{(-8, True)}.

\newpage
\titlepagepart{04 / PYTHON-WERKZEUGE}{Von der Zahl zum formatierten Text}
\textbf{Keine Imports nötig.} \code{bin}, \code{hex}, \code{print}, \code{int},
\code{str} und \code{ValueError} sind eingebaut. Schlüsselwörter sind im Code fett,
Kommentare kursiv; die Darstellung bleibt vollständig schwarz-weiß.

\block{Zahl und Text unterscheiden}
\begin{lstlisting}
value = 6            # int: eine Zahl
text = bin(value)    # str: "0b110"
print(6 + 1)         # 7
print("6" + "1")     # 61: Texte werden aneinandergehaengt
\end{lstlisting}
\code{bin()} und \code{hex()} liefern Text. Die Präfixe \code{0b} und \code{0x}
kennzeichnen die Basis. Sie sind keine zusätzlichen Ziffern des Zahlenwerts.

\block{Die Umwandlung in einzelnen Schritten}
\begin{lstlisting}
mit_praefix = bin(9)               # "0b1001"
ohne_praefix = mit_praefix[2:]      # "1001"
ergebnis = ohne_praefix.zfill(8)   # "00001001"

mit_praefix = hex(15)              # "0xf", kleines f
\end{lstlisting}
\textbf{Indexierung:} Im Text \code{0b1001} haben die Zeichen die Positionen
$0,1,2,3,4,5$. \code{[2:]} übernimmt ab Position 2 bis zum Ende.
Es entfernt das Präfix und rechnet den Text nicht in Dezimal um.
\begin{lstlisting}
ohne_praefix = mit_praefix[2:]      # "f"
gross = ohne_praefix.upper()       # "F"
ergebnis = gross.zfill(2)          # "0F", nicht "FF"
\end{lstlisting}
\code{zfill(n)} ergänzt links Nullen bis zur \textbf{Gesamtlänge} $n$.
\code{"A".zfill(3)} ergibt \code{"00A"}. Es kürzt zu lange Texte nicht und fügt
kein Präfix hinzu. Deshalb prüfen wir den Wertebereich separat.

\block{Wie viele Hexstellen? Erst addieren, dann teilen}
\[
\text{Hexstellen}=\left\lceil\frac{N}{4}\right\rceil
\qquad\text{Python: }\code{(bits + 3) // 4}
\]
\code{//} zählt bei positiven Zahlen die vollständigen Gruppen:
$19=4\cdot4+3$, also \code{19 // 4 == 4}. Der Rest wird nicht aufgerundet.
Das vorherige $+3$ sorgt dafür, dass jeder ursprüngliche Rest von 1, 2 oder 3
noch eine Gruppe mitzählen lässt; bei Rest 0 entsteht keine zusätzliche Gruppe.
\begin{center}
\begin{tabular}{rlll}\toprule
Bits&Vierergruppen&Formel&Hexstellen\\\midrule
8&2 volle&$11\mathbin{//}4$&2\\
9&2 volle, 1 Bit übrig&$12\mathbin{//}4$&3\\
12&3 volle&$15\mathbin{//}4$&3\\
16&4 volle&$19\mathbin{//}4$&4\\\bottomrule
\end{tabular}
\end{center}
\textbf{Klammern!} \code{(8 + 3) // 4} ergibt 2;
\code{8 + 3 // 4} ergibt 8, weil zuerst dividiert wird.

\textbf{Funktionen:} \code{=} speichert einen Wert, \code{==} vergleicht.
\code{return} gibt zurück und beendet den Aufruf; \code{print} zeigt etwas an.
\code{return} gehört in eine Funktion. Einrückung bestimmt die Zugehörigkeit.
Typhinweise wie \code{: int} oder \code{-> str} erzwingen keine Typprüfung.

\newpage
\titlepagepart{05 / REFERENZCODE}{Die ersten drei Funktionen}
Für alle Eingaben setzen wir ganze Zahlen voraus. \code{bits} muss mindestens 1
sein; ein unsigned Muster muss zwischen $0$ und $2^{\mathrm{bits}}-1$ liegen.
\begin{lstlisting}
def to_bin(value: int, bits: int) -> str:
    if bits < 1:
        raise ValueError("Bitbreite muss mindestens 1 sein.")
    if value < 0 or value > 2**bits - 1:
        raise ValueError("Zahl passt nicht in die Bitbreite.")

    mit_praefix = bin(value)
    ohne_praefix = mit_praefix[2:]
    ergebnis = ohne_praefix.zfill(bits)
    return ergebnis


def to_hex(value: int, bits: int) -> str:
    if bits < 1:
        raise ValueError("Bitbreite muss mindestens 1 sein.")
    if value < 0 or value > 2**bits - 1:
        raise ValueError("Zahl passt nicht in die Bitbreite.")

    mit_praefix = hex(value)
    ohne_praefix = mit_praefix[2:]
    gross = ohne_praefix.upper()
    hex_stellen = (bits + 3) // 4
    return gross.zfill(hex_stellen)


def from_twos_complement(raw: int, bits: int) -> int:
    if bits < 1:
        raise ValueError("Bitbreite muss mindestens 1 sein.")
    if raw < 0 or raw > 2**bits - 1:
        raise ValueError("Bitmuster passt nicht in die Bitbreite.")

    grenze = 2**(bits - 1)
    if raw < grenze:
        return raw
    return raw - 2**bits
\end{lstlisting}
\block{Warum wird beim Dekodieren $2^N$ abgezogen?}
Das höchste Bit wechselt seinen Stellenwert von $+2^{N-1}$ zu $-2^{N-1}$.
Die Differenz ist $-2^N$. Alle anderen Stellen bleiben unverändert.
\[
\operatorname{signed}(r,N)=
\begin{cases}
r,&r<2^{N-1},\\
r-2^N,&r\geq2^{N-1}.
\end{cases}
\]
\textbf{Genau an der Grenze beginnt der negative Bereich:} \code{<}, nicht
\code{<=}. Bei vier Bits muss \code{raw=8} zu $-8$ werden.
\begin{lstlisting}
assert to_bin(6, 4) == "0110"
assert to_hex(15, 8) == "0F"
assert from_twos_complement(13, 4) == -3
assert from_twos_complement(13, 8) == 13
\end{lstlisting}

\newpage
\titlepagepart{06 / REFERENZCODE UND PRÜFUNG}{Addition und dein Tagesabschluss}
\begin{lstlisting}
def add_with_overflow_flag(a: int, b: int,
                           bits: int) -> tuple[int, bool]:
    if bits < 1:
        raise ValueError("Bitbreite muss mindestens 1 sein.")

    minimum = -(2**(bits - 1))
    maximum = 2**(bits - 1) - 1
    if a < minimum or a > maximum:
        raise ValueError("a passt nicht in die Bitbreite.")
    if b < minimum or b > maximum:
        raise ValueError("b passt nicht in die Bitbreite.")

    summe = a + b
    overflow = summe < minimum or summe > maximum
    raw = summe % (2**bits)
    ergebnis = from_twos_complement(raw, bits)
    return ergebnis, overflow


assert add_with_overflow_flag(2, 3, 4) == (5, False)
assert add_with_overflow_flag(2, -5, 4) == (-3, False)
assert add_with_overflow_flag(7, 1, 4) == (-8, True)
assert add_with_overflow_flag(-8, -1, 4) == (7, True)
\end{lstlisting}
Die Rückgabe ist ein \textbf{Tupel}: gespeicherte signed Zahl und ein Wahrheitswert.
\code{True} bedeutet „Overflow“, \code{False} bedeutet „kein Overflow“.
Ungültige \emph{Operanden} lösen \code{ValueError} aus; ein Überlauf der
\emph{Summe gültiger Operanden} wird mit dem Flag gemeldet.

\block{Im Terminal ausprobieren}
Aktiviere die Umgebung und wechsle in den Übungsordner. Die folgenden Befehle
sind Shell-Code, kein Python:
\begin{lstlisting}[language=bash,basicstyle=\ttfamily\footnotesize]
source "$HOME/.venvs/ai-hardware-day04/bin/activate"
cd "$HOME/Learning/Semiconductors : AI Hardware/ai-hardware-journey/01_digital/01_zahlensysteme/tag04"
python
\end{lstlisting}
In der Python-Eingabe kannst du dann schreiben:
\begin{lstlisting}
from hardware_numbers.numbers_lernversion import from_twos_complement
print(from_twos_complement(19, 5))  # -13
exit()
\end{lstlisting}
\textbf{Tests:} Im Übungsordner startet \code{python -m pytest -q} die Tests.
Sie prüfen die ursprüngliche Datei \code{numbers.py}, nicht die Lernversion.
Fertige Funktionen dafür in die ursprüngliche Datei übertragen. Lose Experimente
mit undefinierten Variablen dort entfernen.

\block{Was du an Tag 4 beherrschen solltest}
Binär/Hex/Dezimal umrechnen; unsigned und signed Bereiche bestimmen;
Zweierkomplement erklären und dekodieren; addieren, subtrahieren und Carry von
signed Overflow unterscheiden; Vorzeichen korrekt erweitern. Vier Funktionen
implementieren und mindestens 20 Tests inklusive Randfällen bestehen lassen.
Prüfe besonders 0, $-2^{N-1}$, $2^{N-1}-1$, 1 Bit und ungültige Eingaben.

\small\textbf{Grundlage:} Harris \& Harris, \emph{Digital Design and Computer
Architecture, RISC-V Edition}, Abschnitt 1.4, Buchseiten 7--17; dein Lernplan
Woche 1, Tag 4. Eigene Erklärungen und Python-Referenz für unsere Übung.
\end{document}
